\begin{helper}
For every \(s\) and every positive number of colors \(\ell\), the set
\[
  \{n:\noderef{MulticolorRamseyProperty}(s,\ell,n)\}
\]
is nonempty.
\end{helper}
\begin{proof}
We prove the claim by induction on \(\ell\), considering only positive
\(\ell\).  If \(\ell=1\), take \(n=s\).  For any one-color edge-coloring of
\(K_s\), the whole vertex set has size \(s\), and every edge in it has the
unique available color.  Hence
\(\noderef{MulticolorRamseyProperty}(s,1,s)\) holds.

Now suppose the claim is known for a positive number \(\ell\) of colors, and
consider \(\ell+1\) colors.  By the induction hypothesis, choose \(m\) such
that \(\noderef{MulticolorRamseyProperty}(s,\ell,m)\) holds.  By the ordinary
two-color nonemptiness theorem \(\noderef{RamseyPropertyNonempty}\), choose
\(N\) such that every red-blue coloring of \(K_N\) has either a red
\(s\)-clique or a blue \(m\)-clique.

Let \(\chi:E(K_N)\to[\ell+1]\) be any \((\ell+1)\)-coloring.  Build a
two-coloring of the same complete graph by declaring an edge red exactly when
\(\chi\) gives it the last color \(\ell+1\), and blue otherwise.  The
two-color Ramsey property at \(N\) gives two cases.

In the first case there is a red \(s\)-clique.  By the definition of the
two-coloring, every edge of this clique has the last color in \(\chi\), so it
is a monochromatic \(s\)-clique for \(\chi\).

In the second case there is a blue \(m\)-clique \(I\).  On every edge inside
\(I\), the color under \(\chi\) is not the last color, so it can be regarded as
one of the first \(\ell\) colors.  After listing \(I\) as an \(m\)-element
vertex set, this gives an \(\ell\)-coloring of the edges of \(K_m\).  Applying
\(\noderef{MulticolorRamseyProperty}(s,\ell,m)\) to that restricted coloring
produces an \(s\)-element subset whose edges all have some first color
\(c\le\ell\).  Viewed back inside \(K_N\), the same subset is monochromatic
with color \(c\) for the original \((\ell+1)\)-coloring.

Thus every \((\ell+1)\)-coloring of \(K_N\) has a monochromatic \(s\)-clique,
so \(\noderef{MulticolorRamseyProperty}(s,\ell+1,N)\) holds.  This completes
the induction and proves nonemptiness for every positive \(\ell\).
\end{proof}

